\begin{trapbox}
	\begin{description}[style=nextline, leftmargin=2.8em, labelsep=0.5em, itemsep=0.35em]
		\item[\textbf{\textcolor{CTrap}{易错点一（边角对应错误）}}] 在写公式时，将边和角配错，如计算 $a^2$ 时夹角误用了 $\angle B$ 或 $\angle C$。
		\item[\textbf{\textcolor{CTrap}{正确理解（边角对应正确）：}}] 在公式
		      \[
			      a^2=b^2+c^2-2bc\cos A
		      \]
		      中，角 $A$ 是右边两条边 $b,c$ 的夹角，同时也是左边边 $a$ 的对角。也就是说：求哪条边，余弦中就用这条边的对角。
		\item[\textbf{\textcolor{CTrap}{错因分析（记忆模糊）：}}] 对公式结构理解不清，或者未注意公式中的角既是右边两边的夹角，也是左边边的对角。
	\end{description}

	\medskip
	\begin{description}[style=nextline, leftmargin=2.8em, labelsep=0.5em, itemsep=0.35em]
		\item[\textbf{\textcolor{CTrap}{易错点二（忽略符号负号）}}] 在代入余弦时，若夹角为钝角（余弦为负），计算时忘记整体处理公式前面的负号，导致边长平方值出错。
		\item[\textbf{\textcolor{CTrap}{正确理解（符号代入正确）：}}] 代入钝角余弦时，要整体处理公式中的 $-2ab\cos C$。例如若
		      \[
			      \cos C=-\frac12,
		      \]
		      则
		      \[
			      c^2=a^2+b^2-2ab\left(-\frac12\right)=a^2+b^2+ab.
		      \]
		\item[\textbf{\textcolor{CTrap}{错因分析（惯性思维）：}}] 受初中锐角三角函数影响，以为余弦恒正，或只记住 $2ab\cos C$，忽略了公式中的整体负号。
	\end{description}

	\medskip
	\begin{description}[style=nextline, leftmargin=2.8em, labelsep=0.5em, itemsep=0.35em]
		\item[\textbf{\textcolor{CTrap}{易错点三（开方丢失正值）}}] 由 $a^2$ 解 $a$ 时，忘记了 $a>0$，直接给出 $\pm\sqrt{k}$。
		\item[\textbf{\textcolor{CTrap}{正确理解（算术平方根取正）：}}] 若由余弦定理算得
		      \[
			      a^2=k\quad (k>0),
		      \]
		      由于边长为正数，所以
		      \[
			      a=\sqrt{k}.
		      \]
		      一般地，$\sqrt{a^2}=|a|$，在边长问题中因为 $a>0$，所以才取正值。
		\item[\textbf{\textcolor{CTrap}{错因分析（计算习惯）：}}] 习惯了解方程 $x^2 = k$ 得 $x=\pm\sqrt{k}$，忽略了三角形边长的几何约束。
	\end{description}
\end{trapbox}